% Layout helpers for PDF export.  Fonts, margins and colours are supplied by
% the exporter through Pandoc template variables.

\usepackage{etoolbox}
\usepackage{longtable}
\usepackage{booktabs}
\usepackage{array}
\usepackage{pdflscape}
\usepackage{fvextra}
\usepackage{xurl}
\usepackage{graphicx}

% Long code lines and URLs must not run into the margin.
\fvset{breaklines=true,breakanywhere=true}
\DefineVerbatimEnvironment{Highlighting}{Verbatim}{%
  breaklines=true,breakanywhere=true,commandchars=\\\{\}%
}
\setlength{\emergencystretch}{4em}

% Tables: keep cell padding compact, wrap long cells, and use a slightly
% smaller font inside long tables.
\setlength{\tabcolsep}{3pt}
\renewcommand{\arraystretch}{1.15}
\setlength{\LTcapwidth}{\linewidth}
\AtBeginEnvironment{longtable}{\small}

% Images should never overflow the text block.
\setkeys{Gin}{width=\linewidth,height=0.82\textheight,keepaspectratio}

% Arrows and similar symbols are missing from the CJK text fonts; map the
% Unicode characters used in the reports to their math glyphs so that they do
% not render as empty boxes in the PDF export.
\usepackage{newunicodechar}
\newunicodechar{→}{\ensuremath{\rightarrow}}
\newunicodechar{←}{\ensuremath{\leftarrow}}
\newunicodechar{↔}{\ensuremath{\leftrightarrow}}
\newunicodechar{⇒}{\ensuremath{\Rightarrow}}
\newunicodechar{⇔}{\ensuremath{\Leftrightarrow}}

% Enclosed alphanumerics ①-⑳ are used as inline step markers.  They live
% outside xeCJK's default CJK ranges, so they fall back to the Latin font and
% lose their glyph; assigning them to the CJK class keeps them in the CJK font.
\xeCJKDeclareCharClass{CJK}{"2460 -> "2473}

% A remote PDF filename is a Unicode file specification, not a URI or a
% platform byte string. Explicit UTF-16BE /UF avoids ambiguous UTF-8 bytes
% being interpreted as PDFDocEncoding for short Chinese names.
\makeatletter
\newcommand{\ProjectPdfLink}[3]{%
  \leavevmode
  \begingroup
  \Hy@BeginAnnot{%
    \@pdfm@mark{%
      bann<< /Type/Annot /Subtype/Link
        \Hy@setpdfborder
        \Hy@setpdfhighlight
        /A<< /S/GoToR
          /F<< /Type/Filespec /F<#1> /UF<#1> >>
          /D #2
        >>
      >>%
    }%
    \Hy@colorlink\@filecolor
  }%
  #3\Hy@xspace@end
  \Hy@EndAnnot
  \endgroup
}
\makeatother
